% ============================================================
\chapter{Identification, Estimation et Diagnostic}
\label{ch:identification}
% ============================================================

\section{Exemple motivant}
Face \`a une s\'erie temporelle r\'eelle, comment choisir les ordres $p$ et $q$ d'un ARMA$(p,q)$, estimer les param\`etres et v\'erifier que le mod\`ele est ad\'equat ? C'est la \textbf{m\'ethodologie de Box--Jenkins}\index{Box--Jenkins}.

\section{La m\'ethodologie de Box--Jenkins}

\begin{methode}
\begin{enumerate}
\item \textbf{Identification} : examiner l'ACF et la PACF pour proposer des ordres $(p,q)$.
\item \textbf{Estimation} : estimer $(\phi_1,\ldots,\phi_p,\theta_1,\ldots,\theta_q,\sigma^2)$ par maximum de vraisemblance ou moindres carr\'es conditionnels.
\item \textbf{Diagnostic} : v\'erifier que les r\'esidus sont un bruit blanc.
\item \textbf{Pr\'evision} : utiliser le mod\`ele valid\'e pour pr\'evoir.
\end{enumerate}
\end{methode}

\section{Identification}\index{identification}

\subsection{Examen de l'ACF et de la PACF}

R\`egles pratiques (r\'ecapitul\'ees au chapitre~\ref{ch:arma}) :
\begin{itemize}
\item ACF tronque \`a $q$ $\Rightarrow$ MA$(q)$.
\item PACF tronque \`a $p$ $\Rightarrow$ AR$(p)$.
\item Les deux d\'ecroissent $\Rightarrow$ ARMA$(p,q)$.
\end{itemize}

\subsection{Crit\`eres d'information}\index{AIC}\index{BIC}

\begin{definition}[AIC et BIC]
\begin{align}
\AIC(p,q) &= -2\ln L(\hat{\bm\theta}) + 2(p+q+1), \\
\BIC(p,q) &= -2\ln L(\hat{\bm\theta}) + (p+q+1)\ln n.
\end{align}
On choisit le mod\`ele qui \textbf{minimise} le crit\`ere. Le BIC p\'enalise davantage la complexit\'e.
\end{definition}

\begin{theoreme}[Propri\'et\'es asymptotiques]
\begin{itemize}
\item L'AIC est asymptotiquement \'equivalent \`a la validation crois\'ee leave-one-out.
\item Le BIC est \textbf{consistant} : il s\'electionne le vrai mod\`ele avec probabilit\'e $\to 1$.
\item L'AIC tend \`a sur-param\'etrer.
\end{itemize}
\end{theoreme}

\section{Estimation}\index{estimation}

\subsection{Maximum de vraisemblance}\index{maximum de vraisemblance}

\begin{definition}
Pour un ARMA$(p,q)$ gaussien, la log-vraisemblance exacte est :
\[
\ln L(\bm\theta) = -\frac{n}{2}\ln(2\pi) - \frac{1}{2}\sum_{t=1}^n\ln v_t - \frac{1}{2}\sum_{t=1}^n\frac{e_t^2}{v_t},
\]
o\`u $e_t = X_t - \hat{X}_{t|t-1}$ est l'erreur de pr\'evision un pas en avant et $v_t=\Var(e_t)$, calcul\'es r\'ecursivement (algorithme des innovations ou filtre de Kalman).
\end{definition}

\begin{theoreme}[Propri\'et\'es de l'EMV]
Sous des conditions de r\'egularit\'e, l'EMV $\hat{\bm\theta}_n$ est :
\begin{itemize}
\item Consistant : $\hat{\bm\theta}_n\pto\bm\theta_0$.
\item Asymptotiquement normal : $\sqrt{n}(\hat{\bm\theta}_n-\bm\theta_0)\dto\mathcal{N}(0,I(\bm\theta_0)^{-1})$.
\item Asymptotiquement efficace.
\end{itemize}
\end{theoreme}

\subsection{Moindres carr\'es conditionnels}

\begin{definition}
En conditionnant sur les premi\`eres observations, on minimise :
\[
S(\bm\theta) = \sum_{t=p+1}^n\left(X_t - \phi_1 X_{t-1}-\cdots - \phi_p X_{t-p} - \theta_1\hat{\eps}_{t-1}-\cdots-\theta_q\hat{\eps}_{t-q}\right)^2.
\]
\end{definition}

\section{Diagnostic}\index{diagnostic}

\subsection{Analyse des r\'esidus}

\begin{definition}
Les r\'esidus du mod\`ele sont $\hat{\eps}_t = X_t - \hat{X}_{t|t-1}$. Si le mod\`ele est correct, les $\hat{\eps}_t$ doivent se comporter comme un bruit blanc.
\end{definition}

V\'erifications :
\begin{enumerate}
\item \textbf{ACF des r\'esidus} : aucune autocorr\'elation significative.
\item \textbf{Test de Ljung--Box}\index{test!Ljung--Box} :
\[
Q(m) = n(n+2)\sum_{h=1}^m\frac{\hat{\rho}_{\hat\eps}(h)^2}{n-h} \stackrel{H_0}{\sim}\chi^2(m-p-q).
\]
\item \textbf{Normalit\'e} : QQ-plot, test de Jarque--Bera\index{test!Jarque--Bera}.
\item \textbf{H\'et\'erosc\'edasticit\'e} : effet ARCH (cf.\ chapitre~\ref{ch:garch}).
\end{enumerate}

\begin{theoreme}[Ljung--Box]
Sous $H_0$ (r\'esidus = bruit blanc), $Q(m)\sim\chi^2(m-p-q)$ asymptotiquement. On rejette $H_0$ si $Q(m)>\chi^2_{1-\alpha}(m-p-q)$.
\end{theoreme}

\section{Impl\'ementation}

\begin{codeblock}[title=\textbf{Python}]
\begin{minted}{python}
import numpy as np
import pandas as pd
import matplotlib.pyplot as plt
from statsmodels.tsa.arima.model import ARIMA
from statsmodels.stats.diagnostic import acorr_ljungbox
from statsmodels.graphics.tsaplots import plot_acf

# Simuler un ARMA(1,1) connu
np.random.seed(42)
from statsmodels.tsa.arima_process import ArmaProcess
true_ar = np.array([1, -0.7])
true_ma = np.array([1, 0.3])
process = ArmaProcess(true_ar, true_ma)
data = process.generate_sample(nsample=500)

# Étape 1 : Identification (ACF/PACF déjà examinée)

# Étape 2 : Estimation par MLE
model = ARIMA(data, order=(1, 0, 1))
results = model.fit()
print(results.summary())
print(f"\nphi_1 = {results.arparams[0]:.4f} (vrai: 0.7)")
print(f"theta_1 = {results.maparams[0]:.4f} (vrai: 0.3)")

# Étape 3 : Diagnostic
residuals = results.resid

fig, axes = plt.subplots(2, 2, figsize=(12, 8))
axes[0,0].plot(residuals, lw=0.5)
axes[0,0].set_title('Résidus')
plot_acf(residuals, lags=25, ax=axes[0,1], title='ACF des résidus')
axes[1,0].hist(residuals, bins=30, density=True, alpha=0.7)
axes[1,0].set_title('Distribution des résidus')
from scipy import stats
stats.probplot(residuals, dist='norm', plot=axes[1,1])
axes[1,1].set_title('QQ-plot')
plt.tight_layout()
plt.savefig('ch04_diagnostic.pdf')
plt.show()

# Test de Ljung-Box
lb_test = acorr_ljungbox(residuals, lags=20, return_df=True)
print(lb_test)

# Sélection par AIC/BIC
results_table = []
for p in range(4):
    for q in range(4):
        try:
            m = ARIMA(data, order=(p, 0, q)).fit()
            results_table.append({'p': p, 'q': q,
                                  'AIC': m.aic, 'BIC': m.bic})
        except:
            pass
df = pd.DataFrame(results_table).sort_values('BIC')
print(df.head(5))
\end{minted}
\end{codeblock}

\section{Exercices}

\begin{exercice}[$\star$]
Simuler un AR(2) et appliquer la m\'ethodologie compl\`ete : identification, estimation, diagnostic.
\end{exercice}

\begin{exercice}[$\star\star$]
Montrer que le test de Ljung--Box avec $m$ d\'ecalages suit asymptotiquement une $\chi^2(m-p-q)$ sous $H_0$.
\end{exercice}

\begin{exercice}[$\star\star$]
Comparer AIC et BIC sur 1000 simulations d'un AR(1). Lequel s\'electionne le bon ordre le plus souvent ? Le BIC sur-p\'enalise-t-il en petit \'echantillon ?
\end{exercice}

\begin{exercice}[$\star\star\star$ -- Projet]
Appliquer la m\'ethodologie Box--Jenkins compl\`ete \`a la s\'erie du taux de ch\^omage mensuel fran\c{c}ais. Pr\'esenter chaque \'etape avec justification.
\end{exercice}

\begin{formules}
\begin{align*}
\AIC &= -2\ln L + 2k, \quad \BIC = -2\ln L + k\ln n \\
Q(m) &= n(n+2)\sum_{h=1}^m\frac{\hat\rho_{\hat\eps}(h)^2}{n-h}\sim\chi^2(m-p-q)
\end{align*}
\end{formules}
